% ---------- פרק 1: מבוא ורקע ----------
% 22.11.2026: גוף שלישי וקיצור. כל מספר מגובה במקור מ"מקורות/references.tex".
% מבנה משפך: מגמה עולמית ← הארגון ← המוטיבציה ← מבנה הדוח.
\chapter{מבוא ורקע}

\section{המגמה: סיכוני שרשרת האספקה וכניסת הבינה המלאכותית}

שרשראות האספקה של תעשיית התעופה והביטחון נמצאות תחת לחץ גובר: הביקוש למערכות ביטחוניות עולה, והחוסן הפיננסי של רבים מהספקים נשחק. מבחן לחץ של חברת הדירוג \en{RapidRatings} מ-2026 בדק 4,309 ספקים בענף. במצב הנוכחי 17.6\% מהם נמצאים בסיכון פיננסי גבוה או גבוה מאוד, ובתרחיש הלחץ השיעור עולה ל-20.3\%. בחברות הפרטיות השיעור הנוכחי הוא 22.1\%, לעומת 15.9\% בציבוריות, ובתרחיש הלחץ 26.7\% ו-17.9\% בהתאמה \cite{rapidratings2026}. הנתונים מגיעים מחברה שמוכרת דירוגי ספקים ולא ממחקר שפיט, ולא נמצא להם תחליף אקדמי.

הקושי אינו רק בסיכון עצמו, אלא ביכולת לראות אותו. משרד מבקר המדינה האמריקאי (\en{GAO}) מצא ב-2025 שבשרשרת של מטוס ה-\en{F-35} יש בידי משרד ההגנה מידע על פחות מ-10\% מספקי הרכיבים וחומרי הגלם, ובבדיקה נוספת 22 מתוך 63 ספקים שהתבקשו למסור נתונים על מדינת המקור סירבו או לא השיבו \cite{gao2025}. כלומר, גם ארגון בעל משאבים רבים אינו יכול להסתמך רק על מידע שהספק מוסר מרצונו.

ממד נוסף הוא סיכון הסייבר דרך הספקים. סוכנות הסייבר האירופית (\en{ENISA}) ניתחה 24 מתקפות על שרשראות אספקה, ומצאה שבכ-66\% מהמקרים הספקים לא ידעו או לא דיווחו כיצד נפרצו \cite{enisa2021}. ב-2019 עצרה מתקפת כופרה את הייצור במפעלי \en{ASCO Industries}, יצרנית בלגית של חלקי מבנה למטוסים, בכמה מדינות \cite{certeu2019}. גם הרגולציה מגיבה: משרד ההגנה האמריקאי קבע תקן מחייב לאבטחת סייבר אצל ספקיו (\en{CMMC}), שנכנס לתוקף בסוף 2024 \cite{cmmc2024}.

במקביל מתפתחים כלים מבוססי בינה מלאכותית: מודלי שפה גדולים (\en{LLM}) שמזהים סיכונים מתוך כתבות על ספקים \cite{zhao2025}, ומערכות של כמה סוכני \en{AI} שמקשרות אירועים לרשת הספקים של הארגון \cite{almahri2026}. רוב העבודות האלה עוסקות בניטור ספקים קיימים, ולא בהערכת ספק חדש לפני ההתקשרות (פרק~\ref{ch:literature}).

\section{ההקשר היישומי: התעשייה האווירית ותהליך טרום ההתקשרות}

התעשייה האווירית לישראל (תע"א) היא חברה ממשלתית לפי חוק החברות הממשלתיות, ובסוף 2022 עבדה עם כ-5,100 ספקים \cite{mevaker2024}. היא מפתחת ומייצרת מערכות אוויריות, חלליות וביטחוניות. ב-2025 הגיעו הכנסותיה לכ-7.38 מיליארד דולר, כ-66\% מהן מייצוא, וצבר ההזמנות עמד על כ-29 מיליארד דולר \cite{iai2026}.

לפני כל התקשרות עם ספק חדש צריך לענות על שתי שאלות. הראשונה היא \textbf{האם הספק טוב}: האם הוא מייצר מוצרים איכותיים ואמינים, מסוגל לייצר לפי מפרט הרכיב, עומד ביעדים שהציב ויציב כלכלית. השנייה היא \textbf{האם יש סיבה לחשוש ממנו}: סנקציות ומגבלות יצוא, אבטחת הסייבר שלו, יציבות המדינה שבה הוא פועל ויחסיה עם ישראל. מאחר שתע"א עוסקת במערכות להגנת המדינה ובמידע מסווג, ספק לא מתאים עלול לפגוע ברציפות האספקה, באיכות המערכות ובביטחון המידע.

\section{המוטיבציה: למה התהליך הקיים כבר לא מספיק}

בסקר של חברת \en{HICX} מ-2021, בקרב 100 מנהלי רכש בכירים בחברות שמחזורן מיליארד דולר ומעלה, 63\% דיווחו שקליטת ספק חדש נמשכת 15 עד 28 ימים, ו-31\% דיווחו על 29 עד 60 ימים \cite{hicx2021}. הסקר עסק בקליטת ספקים באופן כללי, ולא נמצא נתון פומבי על משך הבדיקה בתעשייה ביטחונית.

הבעיה חריפה במיוחד בחברות פרטיות: שיעור הסיכון הפיננסי בהן גבוה יותר \cite{rapidratings2026}, ודווקא עליהן כמעט אין דוחות כספיים פומביים. \en{Incora}, מפיצה פרטית של חלקים לתעשיית התעופה שסיפקה ליותר מ-8,400 לקוחות, פשטה רגל ב-2023, אחרי ירידה של 18.3\% בהכנסות ונטל חוב כבד \cite{incora2023}.

מכאן הצורך בגישה שתהיה מהירה ושיטתית, תישען על מידע ציבורי שאפשר לאמת, תעריך את \textbf{חוזקות} הספק ולא רק תחפש סיכונים, תזהה \textbf{דגלי סיכון} מכל סוג ותדרג אותם לפי ההשפעה על רציפות האספקה, ותציג במפורש מה לא נמצא ומה צריך לבדוק ידנית. הכלי המוצע בדוח, \textbf{מצפן ספקים}, מפיק הערכה מלאה ומנומקת של הספק, שמובילה להמלצה אם להתקשר איתו.

\section{מבנה הדוח}

פרק~\ref{ch:problem} מגדיר את הבעיה, את שאלת המחקר, את ההשערות ואת מפרט הדרישות. פרק~\ref{ch:literature} סוקר את הספרות ומגדיר את פער הידע. פרק~\ref{ch:market} מציג את השוק ואת המוצרים הקיימים. פרק~\ref{ch:design} בוחר בין חלופות תכן ומתאר את הארכיטקטורה, את מקורות המידע ואת החדשנות. פרק~\ref{ch:workplan} מציג את תוכנית העבודה, פרק~\ref{ch:risks} את הסיכונים, ופרק~\ref{ch:summary} מסכם את המוכנות לשלב המימוש.
